% ===========================================================================
\section{共享信息与剩余不确定性}\label{sec:mutual}
\warmth{0}\quad\whisper{互信息度量一个随机变量包含另一个随机变量的信息量；条件熵度量剩余的不确定性。}

\begin{strand}
互信息 $\MI(X;Y)$ 是一个随机变量包含另一个随机变量信息量的度量，
也就是已知一个随机变量时，另一个随机变量不确定度的减少量。
\end{strand}

\begin{definition}[互信息]\label{def:mi}
设离散随机变量 $X,Y$ 分别取值于 $\X,\Y$，则
\[
  \MI(X;Y)
  =\sum_{x\in\X}\sum_{y\in\Y}
   p_{XY}(x,y)\log\frac{p_{XY}(x,y)}{p_X(x)p_Y(y)}.
\]
等价地，
\[
  \MI(X;Y)=\KL\!\left(p_{XY}\Vert p_Xp_Y\right).
\]
\end{definition}

\block{熵的表达式}
在联合分布下取期望，
\begin{align*}
  \MI(X;Y)
  &=\E\!\left[\log\frac{p_{XY}(X,Y)}{p_X(X)p_Y(Y)}\right]\\
  &=\Ent(X)+\Ent(Y)-\Ent(X,Y).
\end{align*}

对数比中的三个期望完整写为
\[
  \E[\log p_{XY}(X,Y)]
  -\E[\log p_X(X)]
  -\E[\log p_Y(Y)]
  =-\Ent(X,Y)+\Ent(X)+\Ent(Y).
\]

\block{条件熵}

\begin{definition}[条件熵]\label{def:conditional}
$Y$ 在给定 $X$ 时的条件熵为
\[
  \Ent(Y\mid X)
  =-\sum_{x\in\X}\sum_{y\in\Y}
    p_{XY}(x,y)\log p_{Y\mid X}(y\mid x).
\]
等价地，它是各个条件分布之熵的平均：
\[
  \Ent(Y\mid X)
  =\sum_{x\in\X}p_X(x)\Ent(Y\mid X=x).
\]
\end{definition}

\begin{proposition}[链式法则]\label{prop:chain}
\[
  \Ent(Y\mid X)=\Ent(X,Y)-\Ent(X),
  \qquad
  \Ent(X\mid Y)=\Ent(X,Y)-\Ent(Y).
\]
\end{proposition}
\begin{proof}
贝叶斯公式给出
$p_{Y\mid X}(y\mid x)=p_{XY}(x,y)/p_X(x)$。
把这个比值代入定义~\ref{def:conditional}：
\begin{align*}
  \Ent(Y\mid X)
  &=-\sum_{x,y}p_{XY}(x,y)
    \log\frac{p_{XY}(x,y)}{p_X(x)}\\
  &=-\sum_{x,y}p_{XY}(x,y)\log p_{XY}(x,y)
    +\sum_{x,y}p_{XY}(x,y)\log p_X(x)\\
  &=\Ent(X,Y)+\sum_xp_X(x)\log p_X(x)\\
  &=\Ent(X,Y)-\Ent(X).
\end{align*}
交换 $X$ 与 $Y$，同理得到
$\Ent(X\mid Y)=\Ent(X,Y)-\Ent(Y)$。
\end{proof}

\begin{yourturn}
复述链式法则证明中最关键的对数拆分：
\[
  \log\!\bigl(p_X(x)p_{Y\mid X}(y\mid x)\bigr)
  =\answerin[5.2cm]{\log p_X(x)+\log p_{Y\mid X}(y\mid x)}.
\]
\recall{$\log p_X(x)+\log p_{Y\mid X}(y\mid x)$。}
\end{yourturn}

\block{把所有恒等式放在一起}
\begin{strand}
\[
  \MI(X;Y)
  =\Ent(X)+\Ent(Y)-\Ent(X,Y)
  =\Ent(Y)-\Ent(Y\mid X)
  =\Ent(X)-\Ent(X\mid Y).
\]
\end{strand}
